\subsection{Lie derivative and Cartan's magic formula}
In this section we discuss two differential operators on the graded algebra $\Omega^\bullet(M)$ of differential forms on a manifold $M$, namely the Lie derivative $\mathcal{L}_X$ and the interior product $\iota_X$, where $X$ is a smooth vector field on $M$. The main goal of this section is to show well-definedness of $\mathcal{L}_X$ and proof the following theorem, describing the relationships between $\mathcal{L}_X, \iota_X$ and the exterior derivative $d$.

\begin{theorem}
    \label{thm:cartan_calculus}
    Let $M$ be a manifold and $X,Y \in \mathscr{X}(M)$. Then the maps
    \begin{align*}
        d &\colon \Omega^r(M) \to \Omega^{r+1}(M), \\
        \mathcal{L}_X &\colon \Omega^r(M) \to \Omega^r(M), \\
        \iota_X &\colon \Omega^r(M) \to \Omega^{r-1}(M)
    \end{align*}
    satisfy the following Leibniz rules for $\alpha, \beta \in \Omega^\bullet(M)$:
    \begin{align*}
        d(\alpha \wedge \beta) &= d\alpha \wedge \beta + (-1)^{|\alpha|} \alpha \wedge d\beta, \\
        \mathcal{L}_X(\alpha \wedge \beta) &= \mathcal{L}_X\alpha \wedge \beta +  \alpha \wedge \mathcal{L}_X\beta, \\
        \iota_X(\alpha \wedge \beta) &= \iota_X\alpha \wedge \beta + (-1)^{|\alpha|} \alpha \wedge \iota_X\beta,
    \end{align*}
    as well as the following identities:
    \begin{enumerate}[(i)]
        \item $d^2 = 0$, \label{eq:cartan1}
        \item $d\mathcal{L}_X = \mathcal{L}_Xd$,\label{eq:cartan2}
        \item $\mathcal{L}_X = d\iota_X + \iota_Xd$,\label{eq:cartan3}
        \item $\mathcal{L}_X\mathcal{L}_Y - \mathcal{L}_Y\mathcal{L}_X = \mathcal{L}_{[X,Y]}$,\label{eq:cartan4}
        \item $\mathcal{L}_X\iota_Y - \iota_Y\mathcal{L}_X = \iota_{[X,Y]}$,\label{eq:cartan5}
        \item $\iota_X\iota_Y + \iota_Y\iota_X = 0$.\label{eq:cartan6}
    \end{enumerate}
\end{theorem}

Firstly let us restate the definition of the Lie derivative. For a smooth manifold $M$ without boundary, $\alpha \in \Omega^r(M)$ and $X \in \mathscr{X}(M)$ we define
\[
(\mathcal{L}_X \alpha)_p := \frac{d}{dt}\bigg|_{t=0} (\phi_t^\ast \alpha)_p,
\]
where $\phi_t = \phi(t,-)$ denotes the flow of $X$.

\begin{lemma}
    \label{lem:well_definedness_Lie_derivative}
    $(\mathcal{L}_X\alpha)_p$ exists for all $p\in M$ and defines a smooth $r$-form on $M$.
\end{lemma}

Before showing the lemma, let us introduce the notion of a smooth family of $r$-forms.

\begin{definition}[Smooth family of $r$-forms]
    Let $\beta_t$ be a family of $r$-forms on $M$ indexed over some open Interval $I$. We call $\beta_t$ smooth if in some (then all) local coordinates $x^1,\ldots,x^n$ we can write
    \[
    (\beta_t)_p = \sum \beta^{i_1,\ldots,i_r}(t,p)d_px_{i_1} \wedge \ldots \wedge d_px_{i_r}
    \]
    for smooth functions $\beta^{i_1,\ldots,i_r}$. If $I$ is closed we call $\beta_t$ smooth if it can be extended to a smooth family $\beta^\prime_t$ indexed over some open open interval $I^\prime \supseteq I$.
\end{definition}

\noindent Then $t \mapsto (\beta_t)_p$ is a smooth curve in the vector space $\Lambda^r(T_p^\ast M)$, such that 
\[
\left(\frac{d}{dt}\bigg|_{t=t_0} \beta_t\right)_p := \frac{d}{dt}\bigg|_{t=t_0} (\beta_t)_p
\]
is a well-defined $r$-form. Furthermore, in local coordinates we have
\[
    \frac{d}{dt}\bigg|_{t=t_0} \beta_t = \sum \frac{\partial\beta^{i_1,\ldots,i_r}}{\partial t}(t_0,-)dx_{i_1} \wedge \ldots \wedge dx_{i_r},
\]
such that this $r$-form is also smooth and moreover depends smoothly on $t_0$. So 
\[
\dot{\beta_t} := \frac{d}{dt}\beta_t
\]
defines a smooth family of $r$-forms by
\[
\left(\frac{d}{dt}\beta_t\right)_{t_0} := \frac{d}{dt}\bigg|_{t=t_0} \beta_t.
\]

Note that if $\partial M = \emptyset$ we can also directly allow for closed intervals $I$ in the definition above. However we may then always extend the family $\beta_t$ to an open interval (in fact to all of $\R$), as $\beta_t$ can equivalently be considered a smooth $r$-form $\beta$ on $M \times I \subseteq M \times \R$. Under this correspondence we have the following interesting relationship (without proof):
\[
\dot{\beta_t} = \mathcal{L}_{\frac{\partial}{\partial t}}\beta.
\]


\begin{proof}[Proof of Lemma \ref{lem:well_definedness_Lie_derivative}]
    Let $x_1,\ldots,x_n$ be local coordinates on some open neighbourhood $U \subseteq M$ of $p$. Choose a smaller neighbourhood $U^\prime \subseteq U$ of $p$ and $\epsilon > 0$, such that the flow
    \[
    \phi \colon (-\epsilon,\epsilon) \times U^\prime \to U
    \]
    of $X$ is defined. Write
    \[
        \alpha = \sum \alpha^{i_1,\ldots,i_r}dx_{i_1} \wedge \ldots \wedge dx_{i_r}.
    \]
    Then:
    \[
        (\phi_t^\ast\alpha)_p = \sum \alpha^{i_1,\ldots,i_r}(\phi_t(p)) \,d_p(x_{i_1} \circ \phi_t) \wedge \ldots \wedge d_p(x_{i_r} \circ \phi_t)
    \]
    Expanding the differentials 
    \[
        d_p(x_{i_k} \circ \phi_t) = \sum_{j=1}^n \frac{\partial(x_{i_k} \circ \phi_t)}{\partial x_j}(p) d_px_j
    \]
    we see that $\beta_t := \phi_t^\ast \alpha$ is a smooth family of $r$-forms on $U^\prime$ indexed over $t \in (-\epsilon,\epsilon)$. Thus
    \[
    (\mathcal{L}_X \alpha)_p = \left(\frac{d}{dt}\bigg|_{t=0} \beta_t\right)_p
    \]
    exists and depends smoothly on $p$.

\end{proof}

\begin{proposition}[Basic properties of the Lie derivative]
    \label{prop:lie_derivative_properties}
    Let $\alpha \in \Omega^r(M), \beta \in \Omega^s(M)$, $\lambda \in \R$ and $X \in \mathscr{X}(M)$. Then the Lie derivative satisfies the following properties:
    \begin{enumerate}[(i)]
        \item Locality: $(\mathcal{L}_X \alpha)_U = \mathcal{L}_{X|_U} \alpha|_U$ for $U \subseteq M$ open
        \item Linearity in $\alpha$: $\mathcal{L}_X(\lambda \alpha + \beta) = \lambda \mathcal{L}_X \alpha + \mathcal{L}_X \beta$
        \item Leibniz rule: $\mathcal{L}_X(\alpha \wedge \beta) = \mathcal{L}_X \alpha \wedge \beta + \alpha \wedge \mathcal{L}_X \beta$
        \item Compatibility with $d$: $\mathcal{L}_X(d\alpha) = d(\mathcal{L}_X\alpha)$
    \end{enumerate}
\end{proposition}

\begin{proof}
    (i) follows from the fact that locally both sides are given by 
    \[
        \frac{d}{dt}\bigg|_{t=0} \beta_t
    \]
    for the smooth family $\beta_t$ defined in the proof of Lemma \ref{lem:well_definedness_Lie_derivative}. Moreover, (ii) follows immediately from the definition. To see (iii) write
    \[
    \mathcal{L}_X(\alpha \wedge \beta)_p = \frac{d}{dt}\bigg|_{t=0} (\phi_t^\ast\alpha)_p \wedge (\phi_t^\ast\beta)_p.
    \] 
    We claim: For a vector space $V$ and smooth curves $\alpha_t \in \Lambda^r(V^\ast), \beta_t \in \Lambda^s(V^\ast)$ the following holds:
    \[
        \frac{d}{dt}\bigg|_{t=0}(\alpha_t\wedge\beta_t) = \left(\frac{d}{dt}\bigg|_{t=0} \alpha_t\right) \wedge \beta_0 + \alpha_0 \wedge \left(\frac{d}{dt}\bigg|_{t=0} \beta_t\right)
    \]
    W.l.o.g. let 
    \[
    \alpha_t = a(t) \alpha^\prime, \; \beta_t = b(t) \beta^\prime 
    \]
    for some $\alpha^\prime \in \Lambda^r(V^\ast), \beta^\prime \in \Lambda^s(V^\ast)$ and smooth functions $a,b$. Then the general case follows from linearity. We compute:
    \begin{align*}
        \frac{d}{dt}\bigg|_{t=0}(\alpha_t \wedge \beta_t) &= \frac{d}{dt}\bigg|_{t=0} a(t)b(t)\alpha^\prime \wedge \beta^\prime\\
        &= \left( \frac{da}{dt}(0)b(0) + a(0)\frac{db}{dt}(0) \right)\alpha^\prime \wedge \beta^\prime\\
        &=\frac{da}{dt}(0) \alpha^\prime \wedge b(0)\beta^\prime + a(0)\alpha^\prime \wedge \frac{db}{dt}(0)\beta^\prime\\
        &= \left(\frac{d}{dt}\bigg|_{t=0} \alpha_t\right) \wedge \beta_0 + \alpha_0 \wedge \left(\frac{d}{dt}\bigg|_{t=0} \beta_t\right).
    \end{align*}

    This shows the claim and thus (iii). To see (iv) let $\beta_t$ be again the smooth family of $r$-forms defined in the proof of Lemma \ref{lem:well_definedness_Lie_derivative} and write in local coordinates $x_1,\ldots,x_n$:
    \[
    \beta_t = \sum \beta^{i_1,\ldots,i_r}(t,-)dx_{i_1} \wedge \ldots \wedge dx_{i_r}.
    \]
    By linearity we may assume that all functions $\beta^{i_1,\ldots,i_r}$ except for one vanish, i.e.
    \[
    \beta_t = f(t,-)dx_{j_1} \wedge \ldots \wedge dx_{j_r}
    \]
    for some smooth function $f$ and $1 \leq j_1 < \cdots < j_r \leq n$.
    Then
    \[
    d\beta_t = \sum_{k=1}^n \frac{\partial f(t,-)}{\partial x_k} dx^k \wedge dx_{j_1} \wedge \ldots \wedge dx_{j_r}.
    \]
    Thus
    \begin{align*}
        \mathcal{L}_X(d\alpha) &= \frac{d}{dt}\bigg|_{t=0} \phi_t^\ast d\alpha = \frac{d}{dt}\bigg|_{t=0} d\beta_t \\
        &= \sum_{k=1}^n \frac{\partial^2 f}{\partial t\partial x_k}(0,-) dx_k \wedge dx_{j_1} \wedge \ldots \wedge dx_{j_r} \\
        &= d \left(\frac{\partial f}{\partial t}(0,-) dx_{j_1} \wedge \ldots \wedge dx_{j_r}\right) \\
        &= d(\frac{d}{dt}\bigg|_{t=0} \beta_t) = d(\mathcal{L}_X\alpha)
    \end{align*}

\end{proof}

For convenience, let us also restate the definition of $\iota_X$
\begin{definition}
    Let $X \in \mathscr{X}(M)$ and $\alpha \in \Omega^r(M)$. For $r > 0$ define
    \[
    \iota_X \alpha := \alpha(X,-) \in \Omega^{r-1}(M)
    \]
    For $r=0$ let $\iota_X \alpha := 0$.
\end{definition}

Clearly $\iota_X$ is linear and only depends on the local behaviour of $X$ and $\alpha$. Furthermore, it is easily seen that 
\[
\iota_X(\alpha \wedge \beta) = \iota_X\alpha \wedge \beta + (-1)^{\abs \alpha} \alpha \wedge \iota_X\beta.
\]

Thus far the Lie derivative has only been defined on manfiolds without boundary. To extend the definition of $\mathcal{L}_X\alpha$ to manifolds with boundary, first choose extensions $X^\prime$ and $\alpha^\prime$ of $X$ and $\alpha$, respectively, to the double of $M$ and then define
\[
\mathcal{L}_X \alpha := (\mathcal{L}_{X^\prime}\alpha^\prime)|_M
\]
Of course a priori it is not clear that this definition does not depend on the choice of $X^\prime$ and $\alpha^\prime$. However, the following proposition, which is one of the formulas of Theorem \ref{thm:cartan_calculus}, shows that it is in fact independent.

\begin{proposition}[Cartan's magic formula]
    Let $X \in \mathscr{X}(M)$. Then the following identity holds:
    \[\mathcal{L}_X = d\iota_X + \iota_Xd\]
\end{proposition}
\begin{proof}
    Let $\mathcal{D}_X := d\iota_X + \iota_Xd$. From the discussion above it is easy to verify:
    \[
    \mathcal{D}_X(\alpha \wedge \beta) = \mathcal{D}_X\alpha \wedge \beta + \alpha \wedge \mathcal{D}_X\beta.
    \]
    Moreover, $\mathcal{D}_X$ clearly satisfies locality, linearity and commutes with $d$. Thus we can establish 
    \[
    \mathcal{D}_X = \mathcal{L}_X
    \]
    by showing equality on $\Omega^0(M)$. So let $f \in \Omega^0(M)$, then 
    \[
    \mathcal{D}_X f = \iota_Xdf = X \cdot f
    \]
    and
    \[
    \mathcal{L}_Xf = \frac{d}{dt}\bigg|_{t=0} (f \circ \phi_t) = df(X) = X \cdot f
    \]
    This shows the proposition.
\end{proof}

Now let us proof the remaining formulas of Theorem \ref{thm:cartan_calculus}, which are shown similiarly.

\begin{proof}[Proof of Theorem \ref{thm:cartan_calculus}]
    Equation (\ref{eq:cartan1}) is only mentioned for the sake of completeness and will not be shown here. See for example \cite{steimle_geometrie} for a proof. Equations (\ref{eq:cartan2}) and (\ref{eq:cartan3}) have already been shown above. 
    
    To see (\ref{eq:cartan4}) note that again both sides of the equation are differential operators satisfying the same Leibniz rule and commute with $d$, such that it suffices to show equality on $\Omega^0(M)$. But this follows immediately from the definition of the Lie bracket.

    We argue similiarly for (\ref{eq:cartan5}), but in this equation the differential operators on each side do not commute with $d$. However they still satisfy the same Leibniz rule, so we only need to show equality on $\Omega^0(M)$ as well as on the exact one-forms. Clearly both sides vanish on $\Omega^0(M)$ (by definition of the interior product) and for $\alpha = df \in \Omega^1(M)$ we compute:
    \begin{align*}
        (\mathcal{L}_X\iota_Y - \iota_Y\mathcal{L}_X)(df) & = \mathcal{L}_X(Y \cdot f) - (\iota_Y d \mathcal{L}_X)(f) \\
        &= XYf - YXf\\
        &= [X,Y] \cdot f = \iota_{[X,Y]}(df).
    \end{align*}
    Finally, (\ref{eq:cartan6}) is just a way to rewrite
    \[
    \alpha(X,Y,-) = -\alpha(Y,X,-).
    \]
\end{proof}

As an applictaion of Theorem \ref{thm:cartan_calculus} we can show the following identity:

\begin{corollary}
    \label{cartan_cal:cor:formula_d_alpha}
    For $\alpha \in \Omega^1(M)$ and $X,Y \in \mathscr{X}(M)$ the following identity holds:
    \[
    d\alpha(X,Y) = X \cdot \alpha(Y) - Y \cdot \alpha(X) - \alpha([X,Y])
    \]
\end{corollary}
\begin{proof}
    By Theorem \ref{thm:cartan_calculus} (\ref{eq:cartan3}) we have:
    \[
    (\iota_Y\mathcal{L}_X)(\alpha) = Y \cdot \alpha(X) + d\alpha(X,Y).
    \]
    Moreover, applying (\ref{eq:cartan5}) yields:
    \[
    (\iota_Y\mathcal{L}_X)(\alpha) = X \cdot \alpha(Y) - \alpha([X,Y]).
    \]
    The corollary follows.
\end{proof}

\headline{Integrating time dependent forms}

Let $\beta_t$ be a smooth time dependent family of differential forms on a manifold $M$ defined on some interval $I$. At the beginning of the section we described what it means to take the derivative of $\beta_t$ with respect to time. Often one is also interested in the \textit{integral} with respect to time, which is defined similiarly, i.e. pointwise. For $a ,b \in I$ and $p \in M$ we set

\[
\left( \int_{a}^{b} \beta_t dt \right)_p := \int_{a}^{b} (\beta_t)_p dt,
\]
where the right side is just the integral of a smooth curve $[a,b] \to \Lambda^r(T_p^\ast M)$ (see section \ref{subsec:integrating_vector_valued_functions} for a definition of such an integral). Furthermore, in local coordinates we have
\[
    \int_{a}^{b} \beta_t dt = \sum \int_{a}^{b} \beta^{i_1,\ldots,i_r}(t,-) dt \; dx_{i_1} \wedge \ldots \wedge dx_{i_r},
\]
so $\int_{a}^{b} \beta_t dt$ is a smooth $r$-form on $M$.
